\chapter{Polyxan Narrative Generator}
\label{ch:54}

This chapter constructs the Polyxan Narrative Generator: a real-time,
quine-stable narrative engine that converts evolving Polyxan hypergraphs into
coherent story structures.  
Where Chapter~\ref{ch:53} provided the real-time synthetic media pipeline, the
present chapter defines the symbolic and structural substrate that drives
narrative, dialogue, plot progression, and conceptual framing.

The Polyxan Narrative Generator is not a scriptwriter layered atop the system.
It is the \emph{narrative curvature flow} of the Polyxan hypergraph itself.
Every story element emerges from curvature decreases, quine-normalisation, and
RSVP alignment.


\section{Narrative as Curvature Flow}
\label{sec:54:narrative-as-curvature-flow}

The Polyxan hypergraph $\mathcal{G}_t$ at time $t$ encodes the semantic
structure of the current state of meaning.  
Narrative is the temporal ordering of its curvature-driven changes.

\begin{definition}[Narrative State]
A narrative state at time $t$ is
\[
\mathsf{N}_t := (\mathcal{G}_t, \Sigma_t, \kappa_t),
\]
where $\mathcal{G}_t$ is the current hypergraph,
$\Sigma_t$ is its quine-stable core, and
$\kappa_t$ is its discrete curvature distribution.
\end{definition}

The Polyxan Narrative Generator produces narrative transitions as curvature
descends over time.


\section{Narrative Atoms and Quine Frames}
\label{sec:54:atoms}

Narrative atoms are the minimal hypergraph motifs from which all narrative
structure is assembled.

\begin{definition}[Narrative Atom]
A narrative atom is a typed 3-vertex, 2-hyperedge subgraph
\[
a = (v_1,v_2,v_3; e_{12}, e_{23})
\]
whose curvature satisfies
\[
\kappa(a) = 1.
\]
\end{definition}

Narrative atoms correspond to elementary textual units:
relations, actions, implications, percepts, or logical steps.

\begin{definition}[Quine Frame]
A quine frame is a minimal quine-stable subgraph
\[
Q \subseteq \mathcal{G}_t
\]
satisfying $\mathcal{Q}(Q)=Q$ and $\kappa(Q)=0$.
\end{definition}

Quine frames serve as persistent narrative anchors: characters, settings,
topics, and invariants.


\section{Curvature-Driven Plot Progression}
\label{sec:54:plot-progression}

Plot is curvature redistribution.

Each EHR rewrite step (Chapter~\ref{ch:42}) reduces discrete curvature.  
Narrative progression occurs when curvature moves from one motif to another.

\begin{definition}[Plot Step]
A plot step is an EHR rewrite step whose curvature decrease is localised to
a single narrative atom:
\[
\Delta \kappa = -1, \qquad \text{support} = a.
\]
\end{definition}

\begin{theorem}[Curvature Conservation]
The total curvature contributing to narrative flow is constant across EHR
rewrite sequences, except during global quine reduction.
\end{theorem}

Thus plot evolves but the narrative fabric remains coherent.


\section{Narrative Lifts from RSVP Fields}
\label{sec:54:lifts}

The embedding $\iota_t$ (Chapter~\ref{ch:31}) maps hypergraph cells into the
RSVP manifold.  
Narrative is lifted into the continuous layer through these embeddings.

\begin{definition}[Narrative Lift]
Given a narrative atom $a$, its lift is the stratified region
\[
L(a) := \iota_t(\mathrm{Inc}(a)) \subset \mathcal{M}.
\]
\end{definition}

The curvature of $L(a)$ modulates the intensity and pace of the story element
corresponding to atom $a$.

\begin{theorem}[Harmonic Narrative Weight]
The narrative weight of an atom is
\[
w(a) = \int_{L(a)} \mathcal{R}(x) \, dV_x,
\]
and strictly decreases under RSVP flow.
\end{theorem}


\section{Temporal Consistency and Quine-Alignment}
\label{sec:54:temporal-consistency}

Narratives remain interpretable because they remain quine-aligned.

\begin{theorem}[Quine Alignment Theorem]
For any two times $t < t'$, the quine-stable cores satisfy
\[
\Sigma_t \subseteq \Sigma_{t'},
\]
and stabilise to the universal quine $\mathfrak{U}$.
\end{theorem}

This guarantees that narrative threads do not contradict each other as time
progresses.


\section{Narrative Rendering Pipeline}
\label{sec:54:rendering}

The Polyxan Narrative Generator integrates with the media synthesis layer to
produce real-time narrative:

\begin{enumerate}
  \item Curvature analysis determines which atoms are eligible for narrative
        activation.
  \item EHR rewrites select the next active atom(s).
  \item Narrative lifts determine audiovisual expression.
  \item Quine frames ensure interpretability, stability, and continuity.
\end{enumerate}

\begin{definition}[Narrative Frame]
A narrative frame at time $t$ is the tuple
\[
F_t := (\Sigma_t, a_t, L(a_t), w(a_t)).
\]
\end{definition}

Each $F_t$ corresponds to a micro-episode of meaning in the synthetic media
stream.


\section{Long-Form Narrative Convergence}
\label{sec:54:long-form}

Given infinite time, any continuous narrative must re-express the universal
quine in increasingly low-curvature form.

\begin{theorem}[Narrative Convergence]
Any narrative generated by the Polyxan system converges in the limit to a
stationary narrative expressing only the universal quine $\mathfrak{U}$.
\end{theorem}

This theorem is the narrative analogue of the Grand Irreversibility Theorem.


\section{Conclusion}
\label{sec:54:conclusion}

The Polyxan Narrative Generator provides the symbolic backbone for dynamic,
interpretable story construction.  
Narrative becomes curvature descent; plot becomes hypergraph rewriting; meaning
becomes quine stabilisation.  
Together with the continuous RSVP fields and the real-time media pipeline, the
Narrative Generator forms the semantic core of all synthetic media produced by
the system.

Chapter~\ref{ch:55} now extends this symbolic narrative engine into the
continuous visual domain through the Harmonic Scene Composer.

